\chapter{Études de cas I}\label{ch:fm:case-studies-i}

Le 23 septembre 1998, quatorze banques et maisons de titres acceptent de recapitaliser un hedge fund, Long-Term Capital Management, et cinq jours plus tard
elles y engagent environ 3.6~milliards de dollars, soit 90~\% de la valeur du fonds à ce moment-là. Le fonds avait commencé l'année avec 4.67~milliards de dollars de capital et un bilan
28 fois plus grand~; à la fin d'août il avait perdu près de la moitié, et au moment de la recapitalisation presque tout. Chaque chiffre de ce paragraphe était
le résultat de décisions que la firme avait prises bien auparavant~: quel effet de levier pratiquer, à quelles conditions de financement, dans quelles opérations et aux côtés de qui. Ce chapitre
et le suivant étudient trois défaillances de ce type à partir du dossier public, et rejouent chacune en modifiant une seule décision.

\section{Comment lire une défaillance}

Vue de l'extérieur, une défaillance est un récit, et les récits favorisent les coupables désignés et la malchance. Vue de l'intérieur d'une firme, c'est une suite de décisions, dont chacune
paraissait raisonnable au vu des informations de son époque~: une limite fixée, une condition de financement acceptée, une position laissée grandir, un choix de vendre ou
de conserver. La lecture utile s'en tient à ce que le dossier public établit, reconstruit les chiffres à partir de lui, et demande quelle décision, modifiée, aurait
changé l'issue, et de combien. Elle ne rejuge pas ce que les régulateurs et les tribunaux ont constaté, et elle n'attribue pas de mobiles que le dossier
ne mentionne pas.

Deux types de liquidité traversent les trois cas, et la spirale de liquidité du livre~1 est le mécanisme qui les relie.

\begin{definition}[Risque de liquidité de financement, risque de liquidité de marché]\label{def:fm:case-studies-i:liquidity}
Le \emph{risque de liquidité de financement}\index{risque de liquidité de financement} est le risque qu'une firme ne puisse pas honorer ses paiements et ses appels de marge lorsqu'ils arrivent à échéance,
parce que ses prêteurs relèvent les décotes, raccourcissent les échéances ou se retirent, ou parce que ses propres pertes consomment sa trésorerie. \emph{risque de liquidité
de marché}\index{risque de liquidité de marché} est le risque qu'une firme ne puisse pas vendre ni couvrir rapidement une position sans faire évoluer son prix contre elle-même, parce que
la position est importante au regard du volume de marché, ou parce que d'autres détenteurs vendent au même moment.
\end{definition}

Une firme dont le financement est abondant peut attendre que le marché se redresse~; une firme dont les positions peuvent être vendues d'un coup a peu besoin de patience sur son financement. Les défaillances
qui mettent fin aux firmes combinent les deux~: les pertes déclenchent des appels de marge, les appels forcent des ventes, les ventes font bouger les prix, et les nouveaux prix déclenchent d'autres appels
(\cref{fig:fm:case-studies-i:spiral}).

\begin{omfigure}
\begin{tikzpicture}[font=\small,box/.style={draw=omDef,fill=omDef!8,rounded corners,align=center,minimum width=3.3cm,minimum height=0.95cm},
  arr/.style={-{Stealth[length=2.2mm]},thick,omDef}]
\node[box] (loss) at (0,2.2) {pertes sur les\\positions};
\node[box] (call) at (6.4,2.2) {appels de marge,\\décotes relevées};
\node[box] (sell) at (6.4,0) {ventes forcées};
\node[box] (price) at (0,0) {les prix évoluent\\contre les détenteurs};
\draw[arr] (loss) -- node[above,font=\scriptsize] {liquidité de financement} (call);
\draw[arr] (call) -- (sell);
\draw[arr] (sell) -- node[below,font=\scriptsize] {liquidité de marché} (price);
\draw[arr] (price) -- (loss);
\node[draw=omThm,dashed,rounded corners,align=center,font=\scriptsize,text width=3.1cm] at (10.5,1.1) {les autres détenteurs des\\mêmes positions les valorisent\\aux nouveaux prix};
\draw[-{Stealth[length=2mm]},omThm,dashed] (8.85,1.1) -- (8.15,0.35);
\end{tikzpicture}\par
\omcaption{La boucle qui relie les deux types de risque de liquidité~: les pertes font monter les appels de marge, les appels forcent les ventes, les ventes font bouger les prix, et les prix ainsi déplacés sont de nouvelles
pertes, pour le vendeur et pour tous ceux qui détiennent les mêmes positions.}\label{fig:fm:case-studies-i:spiral}
\end{omfigure}

\begin{definition}[Jours de liquidation]\label{def:fm:case-studies-i:days}
Le nombre de \emph{jours de liquidation}\index{jours de liquidation} d'une position est sa taille divisée par le volume que la firme peut échanger chaque jour sans un
impact de marché inacceptable, correspondant en général à une part déclarée du volume quotidien moyen~:
\[ D=\frac{Q}{\pi\,V}, \]
pour une position de $Q$ unités, un volume quotidien moyen de $V$ unités et un taux de participation $\pi$ (celui du livre~10).
\end{definition}

Le nombre de \omterm{def:fm:case-studies-i:days}{jours de liquidation} fait le lien entre la taille d'une position et l'horizon du modèle de risque. Une valeur en risque (VaR) à un jour (livre~6) suppose que la position peut être
sortie dans la journée à des prix proches de ceux du modèle~; une position de vingt \omterm{def:fm:case-studies-i:days}{jours de liquidation} porte vingt jours de risque de prix, et ses propres ventes font bouger le
prix à chacun d'eux.

\section{1998~: effet de levier et conditions de financement}

Le dossier public est ici le rapport du General Accounting Office (aujourd'hui Government Accountability Office) d'octobre 1999. À la fin de 1997, le fonds
restitue environ 2.7~milliards de dollars de capital à ses investisseurs, ce qui laisse des actifs nets de 4.67~milliards de dollars~; son effet de levier au bilan, de 17 pour 1 à la fin de 1994,
atteignait 28 pour 1. Le 31 août 1998, il détenait aussi des dérivés pour environ 1.4~billions de dollars de notionnel hors bilan. De janvier à septembre 1998, il
a perdu près de 90~\% de son capital~: 44~\% en août seul, après le moratoire russe sur la dette de la mi-août, qui a poussé les investisseurs vers la qualité et la liquidité, et la plus grande partie de ce qui restait en septembre.

\begin{dated}{2026-09}{Le dossier public des cas du chapitre}\label{dat:fm:case-studies-i:record}
\textbf{Long-Term Capital Management} (GAO/GGD-00-3, octobre 1999)~: actifs nets de 4.67~milliards de dollars à la fin de 1997, après la restitution d'environ 2.7~milliards de dollars
aux investisseurs~; effet de levier au bilan de 28 pour 1 à la fin de 1997~; 2.3~milliards de dollars d'actifs nets le 31 août 1998, après une perte de 44~\% en août~; le 23
septembre 1998, quatorze banques et maisons de titres ont accepté de le recapitaliser, et le 28 septembre elles ont apporté environ 3.6~milliards de dollars, soit 90~\% de sa
valeur liquidative. \textbf{Amaranth Advisors} (US Senate Permanent Subcommittee on Investigations, hearing record S.~Hrg.~110-235, 2007)~: à certains moments de l'été
2006, il détenait environ 40~\% de tous les contrats NYMEX de gaz naturel en circulation pour l'hiver 2006--2007, 75~\% de novembre 2006 et 60~\% de janvier et
mars 2007, et à certains moments 100\,000 contrats dans un mois~; de la fin août à la mi-septembre 2006, ses positions sur le gaz naturel ont perdu plus de
2~milliards de dollars, et il a liquidé son portefeuille de gaz naturel. \textbf{Global Equity Opportunities} (Goldman Sachs, Form 8-K, 13 août 2007)~: un fonds quantitatif
long/short actions, d'une valeur liquidative d'environ 3.6~milliards de dollars, a reçu un investissement en actions de 3~milliards de dollars de la part de Goldman Sachs et d'autres investisseurs.
\end{dated}

Un bilan égal à 28 fois le capital signifie qu'une baisse de $1/28=3.6~\%$ de la valeur des actifs, nets des couvertures, efface le capital. Les positions du fonds
étaient pour l'essentiel des spread trades, longues d'une obligation ou d'un swap et courtes d'un titre proche, dont la valeur nette bouge bien moins que l'une ou l'autre jambe~; l'effet de levier
était calibré sur ce faible risque net. Ce que mesure ce ratio, c'est la faible marge de manœuvre que laisse le financement lorsque les spreads bougent ensemble et fortement.

Le rapport du GAO explique pourquoi la Federal Reserve a facilité une recapitalisation privée~: elle a jugé qu'une liquidation rapide des positions du fonds
et des positions liées détenues par d'autres acteurs de marché, pourraient faire peser une menace importante sur des marchés déjà fragilisés. Les
créanciers détenaient les mêmes positions, ou en étaient les contreparties~; une vente forcée du fonds aurait aussi déprécié leurs propres positions.

\section{Tutoriel~: trois reconstructions}\label{tut:fm:case-studies-i:tut}

\begin{tutorial}
\textbf{Objectif.} Reconstruire à partir des chiffres publics le capital du fonds de 1998 au fil de l'année et le rejouer à la moitié de l'effet de levier~; reconstruire la
part du fonds de gaz naturel dans les positions ouvertes et ses \omterm{def:fm:case-studies-i:days}{jours de liquidation}~; rejouer le débouclage encombré du livre~7 pour un fonds qui vend et un fonds qui conserve.
\textbf{État final~:} les \cref{fig:fm:case-studies-i:ltcm,fig:fm:case-studies-i:amaranth,fig:fm:case-studies-i:august}.
\begin{tutsteps}
\item \textbf{Capital par période.} À partir de 2.3~milliards de dollars à la fin d'août, après une perte de 44~\%, le capital à la fin de juillet était de $2.3/0.56=\$4.11$
milliards de dollars. Les 3.6~milliards de dollars représentaient 90~\% de la valeur liquidative du 28 septembre, donc la part des investisseurs existants était de $3.6/0.9-3.6=0.40$ milliard de dollars.
\item \textbf{Rendements sur le capital.} De janvier à juillet $4.11/4.67-1=-12.1~\%$~; août $-44~\%$~; de septembre jusqu'au 28 $0.40/2.3-1=-82.6~\%$.
\item \textbf{Effet de levier réduit de moitié.} Si les rendements des actifs restent inchangés, diviser l'effet de levier par deux divise par deux le rendement sur capital de chaque période~:
$4.67\times(1-0.061)(1-0.22)(1-0.413)=2.01$ milliards de dollars le 28 septembre (\texttt{firm.casebook.capital\_path}).
\item \textbf{Concentration.} Les parts de la sous-commission dans les positions ouvertes, et les \omterm{def:fm:case-studies-i:days}{jours de liquidation} pour 100\,000 contrats à 20~\% d'un volume quotidien moyen illustratif
de 25\,000 contrats sur le contrat du mois.
\item \textbf{La sortie encombrée.} Trois fonds dont les portefeuilles long-short représentent 30~milliards de dollars en brut, avec un effet de levier de six et un recouvrement de 0.94~; l'un vend la moitié
de son portefeuille sur cinq jours, les autres conservent (\texttt{firm.casebook.quant\_unwind}, sur \texttt{firm.capacity.unwind}).
\end{tutsteps}
\end{tutorial}

\omcode{code/firm/casebook/firm_casebook.py}{35}{56}{Rendements sur le capital dérivés des chiffres publics du capital, trajectoire du capital à effet de levier proportionnellement réduit,
et les jours de liquidation.}\label{lst:fm:case-studies-i:recon}

\begin{omfigure}
\begin{tikzpicture}
\begin{axis}[width=10.5cm,height=5.4cm,ybar,bar width=17pt,ymin=0,ymax=5.4,ylabel={\small capital (milliards de dollars)},xtick={0,1,2,3},
  xticklabels={fin 1997,fin juill.,fin août,28 sept.},xmin=-0.6,xmax=3.6,tick label style={font=\footnotesize},grid=major,grid style={gray!20},
  table/col sep=comma,nodes near coords,nodes near coords style={font=\tiny,/pgf/number format/fixed,/pgf/number format/precision=2},
  legend style={font=\scriptsize,at={(0.5,-0.22)},anchor=north,/tikz/every even column/.append style={column sep=8pt}},legend columns=2,
  legend cell align=left]
\addplot[fill=omDef!60,draw=omDef] table[x=pos,y=actual]{figdata/desk/28-case-studies-i/ltcm.csv};
\addplot[fill=omThm!45,draw=omThm] table[x=pos,y=half]{figdata/desk/28-case-studies-i/ltcm.csv};
\legend{reconstruit à partir du dossier,à effet de levier réduit de moitié}
\end{axis}
\end{tikzpicture}\par
\omcaption{Le capital du fonds en 1998, reconstruit à partir des chiffres du GAO, comparé aux mêmes rendements de ses actifs à effet de levier réduit de moitié~: 0.40~milliard de dollars
du capital des investisseurs qui subsiste le 28 septembre, contre 2.01~milliards de dollars. Données~: \texttt{fm\_cases1.ltcm}.}\label{fig:fm:case-studies-i:ltcm}
\end{omfigure}

La reconstruction ne laisse aux investisseurs que 0.40~milliard de dollars sur leurs 4.67~milliards de dollars, soit une perte de 91.4~\%, avant la recapitalisation~; la valeur
liquidative de 3.81~milliards de dollars publiée pour le 30 septembre inclut les nouveaux capitaux. Avec un levier réduit de moitié, les mêmes rendements d'actifs auraient laissé 2.01~milliards de dollars,
43~\% du capital, et aucun sauvetage n'aurait été nécessaire. Ce contrefactuel est généreux à un égard et sévère à un autre. Il suppose que les actifs auraient chuté autant,
alors qu'une partie de la chute de septembre venait du marché qui jouait contre un fonds connu pour être en détresse, ce qu'un fonds plus petit n'aurait pas été~;
et elle n'accorde aucun crédit au fonds plus petit pour un moindre nombre de contreparties appelant simultanément des marges. Elle ne retire rien non plus aux rendements des années précédentes,
rendements que le levier réduit de moitié aurait divisés par deux.

\begin{remark}[Effet de levier et durée du financement]\label{rem:fm:case-studies-i:term}
L'effet de levier à lui seul ne met pas fin à un fonds~; ce sont les pertes qui le font lorsque ses prêteurs peuvent exiger des appels de marge plus vite que ses positions ne se redressent. Un fonds au même
levier, mais dont le financement est bloqué pendant des mois et dont les conditions de marge ne peuvent changer du jour au lendemain, dispose de temps~; celui dont les prêteurs peuvent relever les décotes chaque jour a l'\omterm{def:fm:treasury-and-funding:horizon}{horizon de survie}
du chapitre~14. Les deux décisions, le niveau de levier et les conditions auxquelles il est pris, n'en font qu'une.
\end{remark}

\section{2006~: la concentration d'un fonds de gaz naturel}

Le dossier de la sous-commission sénatoriale décrit un hedge fund, Amaranth Advisors, qui a accumulé de très grandes positions sur des contrats de gaz naturel~: parfois 100\,000 contrats sur un seul mois, et, à certains moments de l'été 2006, environ 40~\% de l'ensemble
des contrats NYMEX en circulation pour l'hiver 2006--2007 et 75~\% de ceux de novembre 2006. De fin août à la mi-septembre, les positions
sur le gaz naturel ont perdu plus de 2~milliards de dollars, et le fonds a liquidé l'ensemble de son portefeuille de gaz naturel~; la déclaration liminaire du sénateur Coleman l'opposait à 1998~: la
Réserve fédérale n'a pas eu à intervenir.

\begin{omfigure}
\begin{tikzpicture}
\begin{axis}[width=10.5cm,height=4.8cm,xbar,bar width=10pt,xmin=0,xmax=100,xlabel={\small part des contrats NYMEX en circulation (\%)},
  ytick={0,1,2,3},yticklabels={hiver 2006--07,novembre 2006,janv. 2007,mars 2007},y dir=reverse,ymin=-1.0,ymax=3.6,
  tick label style={font=\footnotesize},grid=major,grid style={gray!20},table/col sep=comma,nodes near coords,nodes near coords style={font=\scriptsize}]
\addplot[fill=omDef!55,draw=omDef] table[y=pos,x=share]{figdata/desk/28-case-studies-i/amaranth.csv};
\draw[gray,dashed] (axis cs:50,-1.0) -- (axis cs:50,3.6) node[pos=0.04,right,font=\scriptsize] {moitié};
\end{axis}
\end{tikzpicture}\par
\omcaption{Les plus fortes parts de contrats NYMEX en circulation détenues par le fonds de gaz naturel à certains moments de l'été 2006, d'après le
rapport de la sous-commission sénatoriale. Données~: \texttt{fm\_casebook.AMARANTH}.}\label{fig:fm:case-studies-i:amaranth}
\end{omfigure}

Une position qui représente 75~\% des positions ouvertes sur un contrat ne trouve d'autre marché où vendre que les 25~\% restants~: le fonds est le marché. Les \omterm{def:fm:case-studies-i:days}{jours de
liquidation} illustrent ce point dans les unités qu'utilise un \omterm{def:fm:the-risk-management-function:cro}{comité des risques}. Avec un volume quotidien moyen illustratif de 25\,000 contrats sur le contrat du mois, 100\,000
contrats prennent 20 jours à un taux de participation de 20~\% et 40 jours à 10~\%~; une mesure de risque à un jour pour une telle position décrit un risque que la firme n'aurait pu
sortir en un jour. Les \omterm{def:fm:limits-and-capital-allocation:conc}{limites de concentration} du chapitre~7 existent précisément pour cela~: une limite sur la part d'une position dans les positions ouvertes et sur ses \omterm{def:fm:case-studies-i:days}{jours de liquidation}
s'active avant une limite sur son risque en dollars, parce qu'elle mesure ce que coûterait la sortie de la position.

\begin{proposition}[Participation et durée de sortie]\label{prop:fm:case-studies-i:days}
Pour une position donnée, les \omterm{def:fm:case-studies-i:days}{jours de liquidation} sont inversement proportionnels au taux de participation~; avec une loi d'impact en racine carrée (livre~10), le coût de
sortie à un taux de participation $\pi$ croît comme $\sqrt\pi$, si bien que diviser ce coût par deux impose de diviser la participation par quatre, ce qui quadruple les jours
et la durée pendant laquelle la position est exposée au marché.
\end{proposition}

\begin{proof}
$D=Q/(\pi V)$ est inversement proportionnel à $\pi$. Sous la loi en racine carrée, l'impact par unité négociée à un rythme quotidien $\pi V$ est proportionnel à
$\sigma\sqrt{\pi V/V}=\sigma\sqrt\pi$, donc le coût total de négociation de $Q$ est proportionnel à $Q\sigma\sqrt\pi$. Le diviser par deux exige $\pi\to\pi/4$, et
alors $D\to4D$.
\end{proof}

\section{Août 2007~: la sortie encombrée}

Pendant la semaine du 6 août 2007, plusieurs hedge funds actions quantitatifs long-short ont subi des pertes sans précédent, suivies d'un rebond important. Le 13
août, Goldman Sachs a publié un communiqué selon lequel les stratégies quantitatives subissaient des pressions liées à une hausse des transactions qui se recoupent, à une flambée de la volatilité et à une hausse des corrélations, que sa réponse avait consisté à réduire le risque et l'effet de levier, et qu'elle et
d'autres investisseurs injectaient 3~milliards de dollars de \omterm{def:fm:the-proprietary-market-making-firm:ownfunds}{fonds propres} dans l'un de ses fonds, Global Equity Opportunities, dont la valeur liquidative s'élevait à environ 3.6~milliards de dollars.
L'étude de Khandani et Lo sur cette semaine, dont traitent les livres~7 et~8, fait l'hypothèse que les pertes ont été déclenchées par le débouclage rapide d'une ou plusieurs
de grands portefeuilles quantitatifs actions market neutral, et que ces pertes ont alors exercé une pression sur un ensemble plus large de portefeuilles actions.

\begin{omfigure}
\begin{tikzpicture}
\begin{axis}[width=10.5cm,height=5.4cm,xlabel={\small jour},ylabel={\small P\&L cumulé (\% des fonds propres)},xmin=0,xmax=20,ymin=-17,ymax=1,
  tick label style={font=\footnotesize},grid=major,grid style={gray!20},table/col sep=comma,legend cell align=left,
  legend style={font=\scriptsize,at={(0.97,0.05)},anchor=south east}]
\addplot[omDef,very thick] table[x=day,y=seller]{figdata/desk/28-case-studies-i/august.csv};
\addplot[omThm,thick,dashed] table[x=day,y=holder]{figdata/desk/28-case-studies-i/august.csv};
\draw[gray,dotted] (axis cs:5,-17) -- (axis cs:5,1) node[pos=0.93,right,font=\scriptsize] {fin des ventes};
\legend{le fonds qui vend la moitié de son portefeuille,un fonds qui détient ses positions}
\end{axis}
\end{tikzpicture}\par
\omcaption{Une sortie encombrée synthétique~: trois fonds aux portefeuilles qui se recoupent (recouvrement 0.94) avec un levier de six~; l'un vend la moitié de son portefeuille sur cinq jours.
Les deux perdent environ 5~\% le premier jour~; le fonds qui détient perd davantage au jour 5 ($-15.1~\%$ contre $-13.6~\%$) et moins au jour 20 ($-7.9~\%$ contre $-9.7~\%$),
à mesure que l'impact temporaire s'estompe. Données~: \texttt{fm\_cases1.august}.}\label{fig:fm:case-studies-i:august}
\end{omfigure}

Le modèle est celui du livre~7, et les chiffres sont synthétiques~: ce qui compte est la forme. Les transactions propres du vendeur font bouger les cours des titres qu'il détient, et
le détenteur, dont le portefeuille recoupe le sien, valorise les mêmes mouvements. Pendant les jours de vente, le détenteur perd un peu plus, puisqu'il conserve encore tout
son portefeuille alors que celui du vendeur se réduit. Quand les ventes cessent, la part temporaire de l'impact s'estompe et les cours remontent en partie~: le détenteur récupère
l'essentiel de cette perte, tandis que le vendeur, qui a vendu au plus bas de son propre mouvement, ne récupère que sur la moitié qu'il a gardée. Au jour 20, le vendeur est en baisse de 9.7~\% et le
détenteur de 7.9~\%.

Ce qu'achète le choix du vendeur, c'est une protection contre le scénario où il ne peut pas tenir~: un prêteur qui réclame des marges, une limite de risque qui impose la réduction, un
investisseur qui demande le rachat de ses parts. Un fonds dont le financement tient tout au long de l'épisode, et dont les limites ont été fixées à l'avance, peut choisir~; un fonds qui en est dépourvu est le vendeur
qu'il choisisse ou non. Le recouvrement est ce qui transforme la décision d'une seule firme en perte pour tous~: à un recouvrement de 0.52, le détenteur perd 4.3~\% au jour 20,
et avec un portefeuille sans rapport il ne perd rien, tandis que la perte propre du vendeur change à peine.

\section{Ce que les trois cas ont en commun}

\begin{method}[Lire une défaillance pour en tirer les décisions]\label{met:fm:case-studies-i:read}
\begin{enumerate}
\item S'en tenir aux sources publiques~: rapports officiels, conclusions judiciaires, dépôts réglementaires. Reporter chaque chiffre dans un registre avec sa source.
\item Reconstituer, à partir de ces chiffres, le capital, les positions et le financement de la firme tout au long de l'épisode, en énonçant chaque dérivation.
\item Lister les décisions prises avant l'épisode qui en ont fixé les termes~: levier, durée du financement, concentration, limites, et qui d'autre détenait les mêmes transactions.
\item Refaire la reconstruction en modifiant une seule décision, et énoncer ce que le contrefactuel suppose et ce qu'il ignore.
\item Transformer le résultat en une limite ou un test que la firme peut appliquer à elle-même~: un couplage du levier et de la durée de financement, une limite de \omterm{def:fm:case-studies-i:days}{jours de liquidation}, un
test de résistance à une sortie encombrée.
\end{enumerate}
\end{method}

Les trois cas partagent une même structure. Chaque firme détenait des positions dont le risque semblait faible en marché normal~: des spreads entre instruments apparentés, des calendar
spreads sur une même matière première, un portefeuille actions couvert. Chacune les détenait à une taille et avec un levier qui faisait de leur sortie un événement de marché. Et chacune dépendait
du comportement des autres~: des prêteurs capables de relever les marges, un marché sans personne en face, des fonds aux mêmes transactions. La décision au niveau de la firme
qui se répète n'est pas la transaction, mais le couplage de la taille, du levier et de la durée de financement avec la capacité du marché à absorber la sortie.

\section{Construction~: le recueil de cas, partie 1}\label{bld:fm:case-studies-i:casebook}

\begin{build}
\bfield{Rôle} Le dossier public des études de cas du livre, sous forme de petits tableaux, chaque chiffre étant rattaché à sa ligne du registre, ainsi que les reconstructions et
contrefactuels des chapitres~28 et~29.
\bfield{Interface} \texttt{firm.casebook}~: \texttt{LTCM}, \texttt{AMARANTH}, \texttt{ltcm\_returns}, \texttt{capital\_path},
\texttt{days\_to\_liquidate}, \texttt{quant\_unwind}~; le simulateur de débouclage de \texttt{firm.capacity}.
\bfield{Règles} Seuls les chiffres du registre entrent dans les tableaux~; chaque chiffre dérivé énonce sa dérivation~; les données illustratives sont signalées comme telles.
\bfield{Tests d'acceptation}
\texttt{code/firm/casebook/tests/}~: les rendements dérivés des chiffres du GAO, la trajectoire du capital à levier ajusté, les \omterm{def:fm:case-studies-i:days}{jours de liquidation}, et le vendeur et le détenteur dans une sortie encombrée.
\bfield{Pour aller plus loin} Un contrefactuel sur la durée du financement, à partir de \texttt{firm.treasury.survival\_horizon}~; une spirale des marges pour le fonds de gaz naturel~; le débouclage avec
plusieurs vendeurs.
\end{build}

\begin{omsources}
\item US General Accounting Office, \emph{Long-Term Capital Management: Regulators Need to Focus Greater Attention on Systemic Risk},
GAO/GGD-00-3, octobre 1999.
\item US Senate, Permanent Subcommittee on Investigations, \emph{Excessive Speculation in the Natural Gas Market}, audition, S.~Hrg.~110-235, 2007.
\item The Goldman Sachs Group, Inc., Form 8-K, Exhibit 99.1, 13 août 2007.
\item A.~E. Khandani et A.~W. Lo, «~What happened to the quants in August 2007~?~», 2007.
\end{omsources}

\section{Exercices}

\begin{exercise}[$\star$]\label{exo:fm:case-studies-i:1}
À partir des chiffres du GAO, calculer le capital du fonds de 1998 à la fin de juillet et son rendement sur capital de la fin d'août au 28 septembre.
\end{exercise}

\begin{exercise}[$\star$]\label{exo:fm:case-studies-i:2}
Définir le \omterm{def:fm:case-studies-i:liquidity}{risque de liquidité de financement} et le \omterm{def:fm:case-studies-i:liquidity}{risque de liquidité de marché}, et dire lequel des deux domine dans chacun des trois cas du chapitre.
\end{exercise}

\begin{exercise}[$\star$]\label{exo:fm:case-studies-i:3}
Calculer le nombre de jours nécessaires pour liquider 100\,000 contrats à 10~\% d'un volume quotidien moyen de 25\,000.
\end{exercise}

\begin{exercise}[$\star\star$]\label{exo:fm:case-studies-i:4}
Avec un levier de 28 pour 1, calculer la baisse de la valeur des actifs du bilan, nets de couvertures, qui efface le capital de 4.67~milliards de dollars, et
en déduire la taille du bilan.
\end{exercise}

\begin{exercise}[$\star\star$]\label{exo:fm:case-studies-i:5}
Démontrer la \cref{prop:fm:case-studies-i:days}, puis en déduire ce qu'elle implique pour une position qui représente 75~\% des positions ouvertes.
\end{exercise}

\begin{exercise}[$\star\star$]\label{exo:fm:case-studies-i:6}
Expliquer pourquoi, dans la \cref{fig:fm:case-studies-i:august}, le détenteur perd plus que le vendeur au jour 5 et moins au jour 20.
\end{exercise}

\begin{exercise}[$\star\star\star$]\label{exo:fm:case-studies-i:7}
\emph{Programmation.} Relancer \texttt{quant\_unwind} avec un recouvrement de 0.5 puis de 0. Calculer les pertes du vendeur et du détenteur au jour 20.
\end{exercise}

\begin{exercise}[$\star\star\star$]\label{exo:fm:case-studies-i:8}
\emph{Trouver la faille.} «~La valeur en risque à un jour de la position ne représente que 5~\% de notre capital, elle ne peut donc pas nous porter un coup sévère.~»
\end{exercise}

\section{Problème~: Trois firmes, trois décisions}

\begin{problem}[{Problème du week-end --- trois firmes, trois décisions}]\label{pb:fm:case-studies-i:1}
Un nouveau \omterm{def:fm:the-risk-management-function:cro}{directeur des risques} (CRO) est prié par le conseil d'administration d'indiquer quelles leçons des échecs quantitatifs les plus célèbres s'appliquent à la firme, et ce qu'il faut modifier.

\medskip\noindent\textbf{Partie I --- La lecture.}
\begin{enumerate}
\item Dire à quoi doit se tenir une lecture utile d'un échec, et ce qu'elle doit éviter.
\item Définir le \omterm{def:fm:case-studies-i:liquidity}{risque de liquidité de financement} et le \omterm{def:fm:case-studies-i:liquidity}{risque de liquidité de marché}.
\item Décrire la boucle qui les relie.
\item Définir les \omterm{def:fm:case-studies-i:days}{jours de liquidation}.
\end{enumerate}

\medskip\noindent\textbf{Partie II --- Le dossier public.}
\begin{enumerate}[resume]
\item Résumer le dossier public du fonds de 1998.
\item Expliquer pourquoi la Réserve fédérale a facilité une recapitalisation privée.
\item Résumer le dossier public du fonds de gaz naturel.
\item Dire ce que Goldman Sachs a déclaré le 13 août 2007.
\end{enumerate}

\medskip\noindent\textbf{Partie III --- Les reconstructions.}
\begin{enumerate}[resume]
\item Calculer les rendements sur capital du fonds de 1998 par période.
\item Donner sa trajectoire de capital, réelle et à moitié du levier.
\item Dire ce que suppose et ce qu'ignore le contrefactuel à demi-levier.
\item Donner les \omterm{def:fm:case-studies-i:days}{jours de liquidation} du fonds de gaz naturel à un taux de participation de 20~\% puis de 10~\%.
\item Énoncer et démontrer la \cref{prop:fm:case-studies-i:days}.
\item Donner les pertes du vendeur et du détenteur aux jours 1, 5 et 20, et expliquer le croisement.
\end{enumerate}

\medskip\noindent\textbf{Partie IV --- Les leçons.}
\begin{enumerate}[resume]
\item Dire ce que les trois cas ont en commun.
\item Dire quelles limites de risque ajouter au cadre de la firme.
\item Expliquer comment soumettre à un test de résistance une sortie encombrée dans le portefeuille propre de la firme.
\item Expliquer pourquoi l'effet de levier et la durée du financement ne relèvent que d'une seule et même décision.
\item Énoncer le \emph{résultat nommé}~: sur la base des chiffres publics, le capital du fonds de 1998 à la fin de septembre s'il avait fonctionné à moitié du levier~; les
\omterm{def:fm:case-studies-i:days}{jours de liquidation} du fonds de gaz naturel à une part déclarée du volume~; et la perte du vendeur d'août 2007 comparée à celle du détenteur.
\item Rédiger la recommandation en deux phrases.
\end{enumerate}
\end{problem}

\section{Questions d'entretien}

\begin{interviewq}[$\star$ \iqroles{risque}]\label{iq:fm:case-studies-i:1}
Quelle est la différence entre le \omterm{def:fm:case-studies-i:liquidity}{risque de liquidité de financement} et le \omterm{def:fm:case-studies-i:liquidity}{risque de liquidité de marché}~? Donnez un exemple de chacun.
\end{interviewq}

\begin{interviewq}[$\star$ \iqroles{trader, risque}]\label{iq:fm:case-studies-i:2}
Un fonds fait tourner son bilan à 25 fois son capital. Quelle baisse de ses actifs l'efface, et pourquoi cela pourrait-il surestimer son risque~?
\end{interviewq}

\begin{interviewq}[$\star\star$ \iqroles{risque}]\label{iq:fm:case-studies-i:3}
Un de vos traders détient 40~\% des positions ouvertes sur un contrat à terme. Comment mesurez-vous le risque, et quelle limite fixeriez-vous~?
\end{interviewq}

\begin{interviewq}[$\star\star$ \iqroles{chercheur, trader}]\label{iq:fm:case-studies-i:4}
Votre portefeuille d'arbitrage statistique (stat arb) perd, pour le troisième jour de suite, trois fois son risque quotidien habituel, et vos concurrents semblent
subir la même chose. Vendez-vous ou conservez-vous~?
\end{interviewq}

\begin{interviewq}[$\star\star$ \iqroles{risque, trader}]\label{iq:fm:case-studies-i:5}
Pourquoi les créanciers d'un fonds le recapitaliseraient-ils plutôt que de le laisser faire faillite~?
\end{interviewq}

\begin{interviewq}[$\star\star\star$ \iqroles{chercheur, risque}]\label{iq:fm:case-studies-i:6}
Concevez un test de résistance d'une sortie encombrée pour un portefeuille actions long-short. Quelles données lui sont nécessaires et que la firme n'observe pas~?
\end{interviewq}
